Research on interwar European trade emphasizes the persistence of commercial relationships despite the Great Depression. We distinguish continuity in trading partners from stability in their relative importance. Using League of Nations export and import declarations for 1927--1937, we compare gravity-based predictions from pre-crisis and annually updated trade maps, conditional on current country margins. The pre-crisis map loses substantial predictive power but remains informative beyond distance and borders. Recent maps predict much better. Reallocation occurs overwhelmingly within established relationships, with Germany gaining regional weight. Strong annual continuity thus coexisted with cumulative reordering: Europe's trade hierarchy changed without wholesale replacement of partners.
